\documentclass[a4paper,12pt]{article}
\usepackage[T1]{fontenc}
\usepackage[utf8]{inputenc}
\usepackage[italian]{babel}
\usepackage{amsmath,cases,eqparbox,amsthm}
\usepackage{geometry}
\geometry{a4paper,top=2cm,bottom=2cm,left=0.8cm,right=1.0cm,%
	heightrounded,bindingoffset=5mm}
\usepackage{graphicx, amsmath, siunitx, tabularx}
\graphicspath{ {./images/} }

\title{Verifica di {MAT} }
\author{crever}

\theoremstyle{plain}
\newtheorem{esercizio}{Esercizio}

\usepackage{xpatch}
\makeatletter
\xpatchcmd{\@thm}{\thm@headpunct{.}}{\thm@headpunct{}}{}{}
\makeatother

\renewcommand{\thefootnote}{\fnsymbol{footnote}}
\pagenumbering{gobble}

\begin{document}
%\maketitle

\begin{center}
	{\textbf{Verifica di {MAT}} \verb!IDV! \textbf{- {IST} }}
\end{center}

\small{\textit{\textbf{ISTRUZIONI:} la valutazione della verifica verrà effettuata tramite la <<Griglia di valutazione>> condivisa nel registro
elettronico. Ad ogni esercizio è associato un punteggio massimo e gli esercizi contrassegnati con \emph{\textbf{*}} rappresentano gli obiettivi minimi. Per ogni esercizio, lo studente deve svolgere soltanto una tipologia \emph{\textbf{[A]}}, \emph{\textbf{[B]}}, ..., che indica il suo livello di difficolt{\`a}, barrandola prima della consegna. Al punteggio realizzato con gli esercizi saranno aggiunti 2 punti, che rappresentano la valutazione di base. Il tempo consentito per lo svolgimento della verifica è di 45 minuti e per gli alunni DSA è previsto un tempo aggiuntivo di 15 minuti}}

\begin{center}
	\begin{tabular}{llll}
		\hline
		&  &  & \\
		Cognome .................................... & 
		Nome ....................................... &
		Classe ......... & 
		Data .................\\
		\hline
	\end{tabular}
\end{center}

% \bigskip

% begin template esercizio
%<<SECESR>>
\begin{tabularx}{\textwidth}{ X X }
	\begin{esercizio}
	\end{esercizio}	
	&
	\\
	
	%<<SECTPL>>	
	\emph{\textbf{LVL}{ASR} {CMD}} 
	&		
	\hspace*{\fill} \verb|(punti: PNT)|
	{ESR}
	\\
	% & \\	
	\hline
	& \\	
	%<<SECTPL>>	

\end{tabularx}
%<<SECESR>>
% end template esercizio

\end{document}
